% cuadro C4_1
\begin{table}[htbp]
\centering
\scriptsize
\caption{Estimación de las finanzas públicas 2025--2026}
\label{cua:C4_1}
\begin{tabular}{lrrrrr}
\toprule
\textbf{Concepto} & \textbf{2025 aprob.} & \textbf{2025 estim.} & \textbf{2026} & \textbf{\% PIB 25a} & \textbf{\% PIB 26} \\
\midrule
RFSP & -1 428 348,1 & -1 559 905,4 & -1 587 349,9 & -3,9 & -4,1 \\
Requerimientos financieros extrapresupuestarios & -257 781,7 & -257 781,7 & -193 579,3 & -0,7 & -0,5 \\
Balance presupuestario & -1 170 566,5 & -1 302 123,7 & -1 393 770,6 & -3,2 & -3,6 \\
Ingresos presupuestarios & 8 055 649,4 & 7 924 967,5 & 8 721 057,3 & 22,3 & 22,5 \\
Ingresos petroleros & 1 142 021,5 & 966 864,4 & 1 204 277,7 & 3,2 & 3,1 \\
Ingresos no petroleros & 6 913 627,9 & 6 958 103,0 & 7 516 779,6 & 19,1 & 19,4 \\
Gobierno Federal (no petroleros) & 5 670 839,1 & 5 728 867,9 & 6 215 701,1 & 15,7 & 16,1 \\
Tributarios & 5 296 426,4 & 5 337 702,5 & 5 838 571,0 & 14,6 & 15,1 \\
No tributarios & 374 412,7 & 391 165,4 & 377 130,1 & 1,0 & 1,0 \\
Organismos y empresas & 1 242 788,8 & 1 229 235,2 & 1 301 078,5 & 3,4 & 3,4 \\
Gasto neto pagado & 9 226 215,8 & 9 227 091,2 & 10 114 827,9 & 25,5 & 26,1 \\
Gasto programable pagado & 6 451 831,3 & 6 465 235,1 & 7 015 853,1 & 17,8 & 18,1 \\
Diferimiento de pagos & -75 800,0 & -75 800,0 & -78 855,7 & -0,2 & -0,2 \\
Gasto programable devengado & 6 527 631,3 & 6 541 035,1 & 7 094 708,8 & 18,0 & 18,3 \\
Gasto no programable & 2 774 384,5 & 2 761 856,1 & 3 098 974,9 & 7,7 & 8,0 \\
Costo financiero & 1 388 373,6 & 1 375 947,3 & 1 572 073,3 & 3,8 & 4,1 \\
Participaciones & 1 340 210,9 & 1 350 108,8 & 1 456 045,9 & 3,7 & 3,8 \\
Adefas & 45 800,0 & 35 800,0 & 70 855,7 & 0,1 & 0,2 \\
Superávit económico primario & 218 307,2 & 74 323,6 & 178 802,6 & 0,6 & 0,5 \\
SHRFSP & 18 591 005,5 & 18 903 594,9 & 20 259 590,7 & 51,4 & 52,3 \\
\bottomrule
\end{tabular}
\\[2pt]\begin{minipage}{0.92\textwidth}\scriptsize Fuente: CGPE 2026, Anexo II.6, p. 84. El renglón del primario está rotulado «económico» en 2026 y «presupuestario» en el CGPE 2025: difieren en 500.0 millones, que son el balance no presupuestario. Elaborado por el ITED.\end{minipage}
\end{table}
